\documentclass[10pt,a4paper]{amsart}
\usepackage[margin=19mm]{geometry}
\usepackage{amsmath,amssymb,mathtools}
\usepackage[scr=rsfs,cal=euler]{mathalfa}
\usepackage{xfrac}
\usepackage[unicode,hidelinks,bookmarks=false]{hyperref}
\newcommand{\ie}{i.e.\ }
\newcommand{\cf}{cf.\ }
\newcommand{\resp}{respectively\ }
\newcommand{\Subsection}{\subsection}
\providecommand{\nobreakdash}{\nobreak\hbox{-}}
\setlength{\emergencystretch}{2em}
\multlinegap0pt
\usepackage{amssymb}
\usepackage{mathtools}
\usepackage[shortlabels]{enumitem}
\newtheorem{proposition}{Proposition}
\newcommand{\C}{\mathbb C}
\newcommand{\R}{\mathbb R}
\newcommand{\Span}{\operatorname{span}}
\newcommand{\U}{\mathcal U}
\newcommand{\Uag}{\mathcal U_{\mathrm{ag}}}
\title{On the Positivity of the Products of Positive Primitive Forms}
\author{Yuhang Liu}
\date{Source: arXiv:2608.14862v2 (CC BY 4.0)}
\begin{document}
\maketitle

\noindent\emph{Source and licence.} On the Positivity of the Products of Positive Primitive Forms. Version arXiv:2608.14862v2 (CC BY 4.0), distributed by arXiv under the Creative Commons Attribution 4.0 International licence, http://creativecommons.org/licenses/by/4.0/, as recorded in the arXiv licence metadata for this exact version and shown on the abstract page https://arxiv.org/abs/2608.14862v2. Full licence notice: \texttt{licenses/arxiv-2608.14862-CC-BY-4.0.txt}.\par\smallskip

\noindent\textit{Editorial scope.} Counted unit: the article's proposition prop:explicit-nonag (source0.tex lines 714--727). The article's complete original proof of it is delivered with it (source0.tex lines 729--790, one \texttt{proof} environment of 62 lines). No mathematics is written, repaired or re-derived: the statement and the proof are the article's own lines, in the article's own notation, and no retained line is substituted or re-flowed.  Retained uncounted as the source-local context the proof needs: lines 268--274.  Nothing was omitted from any retained span, so the package carries no omission record.  No retained line contains a citation, so the wrapper carries no bibliography and no bibliography entry.  No retained line contains a figure environment, an image-inclusion command or drawing code, so no image is delivered and the package declares no asset dependency.  Editorial formatting only, in addition: an AMS \texttt{amsart} preamble that restates the article's own declarations and macros used by the retained lines, namely the environments \texttt{proposition} on separate counters, together with the standard packages those lines need.  The retained environments print in this document as Proposition 1 in order of appearance, whereas the article prints the same items with the numbering of its own full document.  The complete native source of the article as distributed in the arXiv e-print for this exact version is bundled unchanged as \texttt{sources/207694/source0.tex}.
\begin{proposition}\label{prop:stable-real-parts}
Let $\dim_{\R}V=6$ and let
$\Xi=\gamma+i\delta\in\U(V)$.  Every nonzero
$\rho\in\Span_{\R}\{\gamma,\delta\}$ is the real part of a nonzero
$(3,0)$-form for an $\omega$-compatible complex structure.  This complex
structure is uniquely determined by $\rho$.
\end{proposition}

\begin{proposition}\label{prop:explicit-nonag}
On $V=\R^6=\C^3$, with
$\omega=\sum_{j=1}^3dx_j\wedge dy_j$ and $z_j=x_j+iy_j$, fix
$0<a<1$ and set
\begin{align*}
 Z&:=dz_1\wedge dz_2\wedge dz_3,\\
 W_a&:=(a\,dx_1+i\,dy_1)\wedge dz_2\wedge dz_3,\\
 \Omega_a&:=\operatorname{Re}Z+i\,\operatorname{Im}W_a.
\end{align*}
Then
\[
 \Omega_a\in\U^+(V)\setminus\Uag(V).
\]
\end{proposition}

\begin{proof}
Both $Z$ and $W_a$ are primitive.  Indeed,
$(a\,dx_1+i\,dy_1)\wedge dx_1\wedge dy_1=0$, while for $j=2,3$ one has
$dz_j\wedge dx_j\wedge dy_j=0$.  Hence $W_a\wedge\omega=0$, and
$\Omega_a$ is primitive.

Write
\[
 T:=d\overline z_1\wedge dz_2\wedge dz_3,
 \qquad
 c:=\frac{1+a}{2},\qquad
 \lambda:=\frac{a-1}{2},
\]
so that $W_a=cZ+\lambda T$.  Let $L\subset\C^3$ be Lagrangian, and write
an oriented orthonormal basis of $L$ as the columns of a unitary matrix
$U$.  Then $Z|_L=\det U\neq0$.  Replacing the first row by its complex
conjugate and applying the row-replacement formula gives
\[
 \frac{T|_L}{Z|_L}
 =\overline{\sum_{j=1}^3U_{1j}^2}=: \zeta_L,
 \qquad |\zeta_L|\leq1.
\]
Consequently,
\begin{equation}\label{eq:Wa-positive}
 \operatorname{Re}\left(\frac{W_a|_L}{Z|_L}\right)
 =\operatorname{Re}(c+\lambda\zeta_L)
 \geq c-|\lambda|=a>0.
\end{equation}
If $\Omega_a|_L=0$, then $Z|_L=ir$ for some nonzero real number $r$, while
\eqref{eq:Wa-positive} gives
\[
 \operatorname{Im}(W_a|_L)
 =r\operatorname{Re}\left(\frac{W_a|_L}{Z|_L}\right)\neq0,
\]
a contradiction.  The same estimate holds while $a$ varies continuously
from its given value to $1$, and $\Omega_1=Z$.  Hence
$\Omega_a\in\U^+(V)$.

The real form $\operatorname{Re}Z$ determines the standard compatible
complex structure $J_1$.  The form
$\operatorname{Im}W_a=\operatorname{Re}(-iW_a)$ determines the compatible
complex structure $J_a$ which is standard on the last two coordinate
planes and satisfies
\[
 J_a\frac{\partial}{\partial x_1}
 =a\frac{\partial}{\partial y_1},
 \qquad
 J_a\frac{\partial}{\partial y_1}
 =-a^{-1}\frac{\partial}{\partial x_1}.
\]
Indeed, $a\,dx_1+i\,dy_1$, $dz_2$, and $dz_3$ are of type $(1,0)$ for
$J_a$.  Thus $J_a\neq J_1$.  If $\Omega_a$ were almost geometric, then
\[
 \Span_{\R}\{\operatorname{Re}Z,\operatorname{Im}W_a\}
 =\Span_{\R}\{\operatorname{Re}\Psi,\operatorname{Im}\Psi\}
\]
for a compatible $(3,0)$-form $\Psi$.  Every nonzero element of the
right-hand plane induces the same complex structure by the uniqueness in
Proposition~\ref{prop:stable-real-parts}.  This contradicts the distinct
complex structures induced by $\operatorname{Re}Z$ and
$\operatorname{Im}W_a$.  Therefore $\Omega_a\notin\Uag(V)$.
\end{proof}

\end{document}
